\section{Methodology: implementation and deployment}
\label{sec:implementation}

This chapter describes how the design is implemented and deployed. For further details on the setup
and on reproducing every step and experiment, see the file
\href{https://github.com/Cloud-Computing-group-10/financial-agent-project/blob/feature/cloud-app/REPRODUCE.md}{\code{REPRODUCE.md}}
in our source code.

\subsection{Tech stack}

The system is written in Python. \Cref{tab:stack} lists its main components. All LLM calls go through
one module, so the same workflow code runs with the hosted Claude model, with Qwen on a laptop, or
with Qwen served on the cloud machine. One Docker image contains the code, Spark and the embedding
model, and the same Docker Compose file starts the system on a laptop and on AWS. Automated tests
check the rules that decide correctness (periods, metric tags, routing, company lookup, passage
splitting, the formula evaluator and the answer check) and run on every change to the code.

\begin{table}[ht]
\centering\small
\begin{tabularx}{\linewidth}{lX}
\toprule
Component & Used for \\
\midrule
Postgres 16 + pgvector & Figures, filings and passages with their embeddings, in one database \\
PySpark (local mode) & The two batch jobs: XBRL figures to tables, and 10-Q text to embedded passages \\
fastembed (MiniLM-L6 on CPU) & Embeddings of passages and questions \\
BeautifulSoup, pypdf & 10-Q web pages and FinanceBench PDFs to text \\
FastAPI, Streamlit & The HTTP service and the web page \\
Docker, Docker Compose & One image for every part; the same setup on the laptop and on EC2 \\
Claude Haiku 5.5 (Anthropic API) & Hosted LLM for the system and for grading \\
Qwen2.5 3B and 7B & Local LLMs: on a MacBook GPU (mlx-lm) or on the EC2 CPU (Ollama) \\
pytest, GitHub Actions & 37 automated tests on every push \\
\bottomrule
\end{tabularx}
\caption{Main components.}
\label{tab:stack}
\end{table}

\subsection{Data pipeline}

One script runs the whole pipeline on the cloud machine (\cref{fig:pipeline,tab:pipeline}). Every
step skips work that is already done, so an interrupted run simply continues.

\begin{figure}[ht]
\centering
\begin{tikzpicture}[
  st/.style={draw,rounded corners=2pt,align=center,font=\footnotesize,text width=26mm,minimum height=11mm},
  sp/.style={st,fill=llmbox},
  arr/.style={-{Stealth[length=2mm]},thick}]
  \node[st] (fetch) at (0,0)     {SEC bulk files};
  \node[sp] (fj)    at (3.6,0)   {Spark facts job\\JSON $\to$ Parquet};
  \node[st] (lf)    at (7.2,0)   {Load figures};
  \node[st] (pg)    at (10.8,-1.2) {Postgres\\+ pgvector};
  \node[st] (man)   at (0,-2.4)  {List of\\140,226 10-Qs};
  \node[st] (dl)    at (3.6,-2.4) {Download\\8 requests/s};
  \node[sp] (tj)    at (7.2,-2.4) {Spark text job\\split + embed};
  \draw[arr] (fetch)--(fj); \draw[arr] (fj)--(lf); \draw[arr] (lf)--(pg);
  \draw[arr] (fetch)--(man); \draw[arr] (man)--(dl); \draw[arr] (dl)--(tj); \draw[arr] (tj)--(pg);
\end{tikzpicture}
\caption{The data pipeline. Shaded steps are Spark jobs. The text job processes downloaded documents
while later ones are still downloading.}
\label{fig:pipeline}
\end{figure}

\begin{table}[ht]
\centering\small
\begin{tabularx}{\linewidth}{Xr}
\toprule
Step & Time on c7i.8xlarge \\
\midrule
Download the SEC bulk files (\code{companyfacts.zip}, \code{submissions.zip}) & under 2 min \\
Spark facts job: 20,609 company files $\to$ Parquet & 40 s \\
Load the figures into Postgres & 60 s \\
List all 10-Qs since 2019 (140,226 filings) & 20 s \\
Download the 10-Qs at 8 requests/s, in 71 batches of 2,000 & about 5 h \\
Spark text job: split into passages, embed, load & about 8 h (overlaps the download) \\
\bottomrule
\end{tabularx}
\caption{Pipeline steps and their run time on the full data.}
\label{tab:pipeline}
\end{table}

The download of the 10-Qs stays just under the SEC's limit of 10 requests per second and stores each
document compressed. Computing embeddings is the most expensive step, so for most companies we keep
only the MD\&A section, which contains the explanations, rather than every section. This keeps about
30\% of the text. For the ten companies of our question set and the FinanceBench filings we keep
every section.

\subsection{Data processing with Spark}
\label{sec:spark}

Both batch jobs apply the same Python function to many independent input files and then combine the
results. Spark divides the files into partitions and processes them in parallel on all cores of the
machine.

\begin{figure}[ht]
\centering
\begin{tikzpicture}[
  st/.style={draw,rounded corners=2pt,align=center,font=\footnotesize,text width=25mm,minimum height=12mm},
  sp/.style={st,fill=llmbox},
  arr/.style={-{Stealth[length=2mm]},thick}]
  \node[st] (in)  at (0,0)     {20,609 JSON files\\(1 per company)};
  \node[sp] (rd)  at (3.3,0)   {\code{wholeTextFiles}\\128 partitions};
  \node[sp] (fm)  at (6.6,0)   {\code{flatMap}\\extract figures};
  \node[sp] (ca)  at (9.9,0)   {\code{cache}\\(table, row) pairs};
  \node[sp] (fi)  at (9.9,-2.2) {\code{filter} by table\\$\to$ DataFrame};
  \node[st] (pq)  at (6.6,-2.2) {Parquet: companies, documents, figures, tags};
  \node[sp] (gb)  at (3.3,-2.2) {\code{groupBy} +\\\code{agg}};
  \node[st] (rep) at (0,-2.2)  {Coverage and restatement reports};
  \draw[arr] (in)--(rd); \draw[arr] (rd)--(fm); \draw[arr] (fm)--(ca); \draw[arr] (ca)--(fi);
  \draw[arr] (fi)--(pq); \draw[arr] (pq)--(gb); \draw[arr] (gb)--(rep);
\end{tikzpicture}
\caption{The Spark facts job. Shaded boxes are Spark operations; each partition is processed by one
core.}
\label{fig:spark}
\end{figure}

\paragraph{Facts job.} The job reads the 20,609 company files and splits them into 128 partitions,
four per core, so that every core stays busy although the files range from a few kilobytes to 30 MB.
A \code{flatMap} turns each file into rows for the four database tables. The result is kept in memory
(\code{cache}) because it is used once for each table, and each table is written as a Parquet file.
The job produces 9,695 companies, 140,436 10-Qs, 3,022,763 figures and 96,950 tag choices in 40
seconds. Two aggregations (\code{groupBy}) then report how many companies report each metric, and how
often a later filing reports a different value for the same quarter (62,483 cases). These
\emph{restatements} are why every figure keeps its filing.

\paragraph{Text job.} The text job reads the downloaded 10-Qs, splits every document into passages
and redistributes them evenly over the partitions, because documents differ greatly in length. Each
partition then embeds its passages in groups of 256, so the embedding model is loaded once per task
instead of once per passage.

\subsection{Updating with new filings}
\label{sec:update}

Companies file new 10-Qs every day. An update script runs the same pipeline on new data only: it adds
the figures of every filing newer than the newest one in the database and downloads and embeds only
the new documents. The facts job still reads all companies, because the SEC offers no file of recent
changes, but it takes under a minute.

\subsection{Deployment on AWS}

\Cref{tab:instance} shows the machine we used. Embedding millions of passages is CPU-bound and
parallel, so we chose a compute-optimised instance with many cores; E5 (\cref{sec:e5}) confirms that
the Spark jobs use up to its 16 physical cores. 64 GB of memory hold Spark and Postgres side by side,
and the 250 GB disk leaves ample room for the data in \cref{tab:storage}. After the full load the
same disk can be attached to a much smaller instance for serving, since answering a question needs
little computation.

\begin{table}[ht]
\centering\small
\begin{tabularx}{\linewidth}{lX}
\toprule
Setting & Value \\
\midrule
Instance type & c7i.8xlarge (compute-optimised): 32 vCPUs on 16 physical cores, 64 GB memory \\
Storage & 250 GB general-purpose SSD (gp3) \\
Operating system & Ubuntu, Docker Compose \\
Region & us-east-1 \\
Network access & SSH and the web page port from our own IP address only; database not exposed \\
Permissions & IAM role with access to our S3 bucket, so no keys are stored on the machine \\
\bottomrule
\end{tabularx}
\caption{EC2 instance.}
\label{tab:instance}
\end{table}

\begin{table}[ht]
\centering\small
\begin{tabularx}{\linewidth}{Xr}
\toprule
Data & Size \\
\midrule
\code{submissions.zip} & 1.5 GB \\
\code{companyfacts.zip}, unpacked into one JSON file per company & 19 GB \\
140,226 downloaded 10-Q documents & 17 GB \\
Passages with embeddings written by Spark (Parquet) & 8.7 GB \\
Postgres database & 12 GB \\
\quad of which passage embeddings (4.55 million $\times$ 384 numbers) & 6.7 GB \\
\quad of which passage text & 3.1 GB \\
\quad of which figures (3.0 million) and the other tables & 0.9 GB \\
Docker images (our image 2.6 GB, Postgres, Ollama with Qwen2.5-3B) and build cache & 22 GB \\
\midrule
Disk in use in total (\code{df}) & 77 GB of 242 GB \\
\bottomrule
\end{tabularx}
\caption{Data stored on the instance, measured with \code{du}, \code{docker system df} and Postgres'
\code{pg\_total\_relation\_size} after the full load and the update run. The database is smaller
than the raw files it was built from; the raw files and Parquet output could be moved to S3 after
loading.}
\label{tab:storage}
\end{table}

The FinanceBench PDFs were processed on a laptop and written into the cloud database through the SSH
tunnel, which the local Qwen models also used to reach the database during the experiments.
